\documentclass{article}

\usepackage[spanish, es-noquoting]{babel}

\usepackage[letterpaper,top=2cm,bottom=2cm,left=3cm,right=3cm,marginparwidth=1.75cm]{geometry}

\title{Arquitectura Tarea (Lista de comandos)}
\author{Jesús Romero  C.I. 32.590.184}

\begin{document}
\maketitle


\begin{enumerate}

    \item add (add)
    \begin{itemize}
        \item Definición: Suma de datos en registros.
        \item Ejemplo: add \$s1,\$s2,\$s3
        \item Significado: \$s1 = \$s2 + \$s3
    \end{itemize}
    
    \item substract (sub)
    \begin{itemize}
        \item Definición: Resta de datos en registro.
        \item Ejemplo: sub \$s1,\$s2,\$s3
        \item Significado: \$s1 = \$s2 - \$s3
    \end{itemize}

    \item add immediate (addi)
    \begin{itemize}
        \item Definición: Usado para sumar constantes.
        \item Ejemplo: addi \$s1,\$s2,100
        \item Significado: \$s1 = \$s2 + 100
    \end{itemize}

    \item load word (lw)
    \begin{itemize}
        \item Definición: Convertir una palabra de memoria a registro.
        \item Ejemplo: lw \$s1,100(\$s2)
        \item Significado: \$s1 = Memory[\$s2 + 100]
    \end{itemize}

    \item store word (sw)
    \begin{itemize}
        \item Definición: Convertir una palabra de registro a memoria.
        \item Ejemplo: sw \$s1,100(\$s2)
        \item Significado: Memory[\$s2 + 100] = \$s1
    \end{itemize}
    
    \item shift left logical (sll)
    \begin{itemize}
        \item Definición: Desplazamiento a la izquierda por constante.
        \item Ejemplo: sll \$s1,\$s2,10
        \item Significado: \$s1 = \$s2 $<<$ 10   
    \end{itemize}

    \item shift right logical (srl)
    \begin{itemize}
        \item Definición: Desplazamiento a la derecha por constante.
        \item Ejemplo: srl \$s1,\$s2,10
        \item Significado: \$s1 = \$s2 $>>$ 10
    \end{itemize}

    \item branch on equal (beq)
    \begin{itemize}
        \item Definición: Comprueba igualdad y bifurca relativo al PC.
        \item Ejemplo: beq \$s1, \$s2, L
        \item Significado: if (\$s1 == \$s2) go to L \\
                           \hspace*{5em} PC + 4 + 100
    \end{itemize}

    \item branch on not equal (bne)
    \begin{itemize}
        \item Definición: Comprueba si no es igual y bifurca relativo al PC.
        \item Ejemplo: \$s1,\$s2,L
        \item Significado: if (\$s1 != \$s2) go to L \\
                           \hspace*{5em} PC + 4 + 100
    \end{itemize}

    \item set on less than (slt)
    \begin{itemize}
        \item Definición: Compara si es menor que un registro.
        \item Ejemplo: slt \$s1,\$s2,\$s3
        \item Significado: if (\$s2 < \$s3) \$s1 = 1;
                           else \$s1 = 0
    \end{itemize}

    \item set on less than immediate (slti)
    \begin{itemize}
        \item Definición: Compara si es menor que una constante
        \item Ejemplo: slt \$s1,\$s2,100
        \item Significado: if (\$s2 < 100) \$s1 = 1;
                           else \$s1 = 0
    \end{itemize}
    
\end{enumerate}

\end{document}
